\section{Uniform concentration of evaluation flats}
\label{conc:section}

We prove the algebraic estimate used twice in the distance argument.  Its
constant is uniform in the scale parameter.  The distinction between the
two versions below is essential: the global family has a cardinality
bound, whereas a family inside one of its flats need not have one.

Throughout this section, varieties are reduced complex algebraic
varieties.  The degree of an affine variety means the degree of its
projective closure; the degree of a pure-dimensional union is the sum of
the degrees of its irreducible components.  Constants denoted by $C_e$
may change from line to line and depend only on $e$.

Let $E$ be a complex vector space of dimension $e\geq3$, and write
\[
 S=\bigwedge^2E^*,\qquad n=\binom e2,\qquad
 r=e-1,\qquad s_0=n-r=\binom{e-1}{2}.
\]
An element $z\in S$ is regarded as a skew map $E\to E^*$.  Set
\[
 \mathcal Q=\{[v:y]\in\mathbf P(E\oplus E^*):y(v)=0\}.
\]
For $\lambda=[v:y]\in\mathcal Q$ with $v\ne0$, define
\begin{equation}
 F_\lambda=\{z\in S:zv=y\},\qquad
 K_v=\{z\in S:zv=0\}.
 \label{conc:flat-model}
\end{equation}
Evaluation at $v$ maps $S$ onto the $r$-dimensional annihilator of $v$.
Consequently $F_\lambda$ is a nonempty affine $s_0$-flat, with direction
$K_v=\bigwedge^2(E/\mathbb Cv)^*$.  Distinct parameters give distinct
flats.  Indeed, the common radical of the forms in $K_v$ is exactly
$\mathbb Cv$: the quotient has dimension at least two, and each of its
nonzero vectors is detected by an alternating form.  The direction of
the flat therefore recovers $[v]$, and its prescribed evaluation recovers
$y$ with the same normalization.

A maximal projective linear subspace of $\mathcal Q$ is a
$\mathbf P^r$.  Equivalently, it is the projectivization of an
$e$-dimensional totally isotropic subspace for the quadratic form
$(v,y)\mapsto y(v)$.

\begin{theorem}[Uniform concentration]
\label{conc:theorem}
Fix $C_0\geq1$.  Let $A\geq1$ and let $\mathcal P$ be a finite set of
distinct parameters $[v:y]\in\mathcal Q$, all with $v\ne0$.  Assume:
\begin{enumerate}
\item[(i)] For every $\lambda\in\mathcal P$ there is a homogeneous polynomial
      of positive degree at most $C_0A$ which is nonzero at $\lambda$ and
      zero at every other member of $\mathcal P$.
\item[(ii)] Every projective $\mathbf P^r\subset\mathcal Q$ contains at most
      $C_0A^{r-1}$ members of $\mathcal P$.
\end{enumerate}
There is a constant $C=C(e,C_0)$ such that
\begin{equation}
 \#\{\lambda\in\mathcal P:F_\lambda\subset V\}
 \leq C\deg(V)A^h,
 \qquad \dim V=s_0+h,\quad 1\leq h<r,
 \label{conc:bound}
\end{equation}
for irreducible proper affine varieties $V\subset S$, in either of these
settings:
\begin{enumerate}
\item[(a)] All such $V$ are allowed, and $|\mathcal P|\leq C_0A^r$.
\item[(b)] Only $V$ contained in the zero set of some nonzero polynomial
      of degree at most $C_0A$ are allowed.  No bound on
      $|\mathcal P|$ is required.
\end{enumerate}
\end{theorem}

For $h=0$, an irreducible variety containing an $s_0$-flat must equal
that flat, so it contains at most one member of the family.  In (a) the
ambient endpoint $h=r$ is supplied by the cardinality assumption.  As a
useful consequence of (b), a nonzero polynomial of degree $D\leq C_0A$
contains at most $CDA^{r-1}$ of the flats in its zero set: apply
\eqref{conc:bound} to its irreducible hypersurface components and sum
their degrees.

The proof first obtains approximate equations and regular defining
equations (Theorem~\ref{aci:theorem} and Lemma~\ref{conc:regularity}),
then recovers controlled families of actual rulings
(Lemma~\ref{conc:rulings}). Proposition~\ref{conc:degree} bounds their degree
and identifies the maximal-isotropic exception.
Section~\ref{conc:uniform-completion} closes the resulting finite
system of inequalities with constants independent of $A$.

We normalize the number of flats contained in a variety by its degree,
following Walsh's concentration parameters
\cite[Section~1.1, definition of \(\mathcal D_m\)]{walsh2023}.  The controlled-ruling strategy and the
isotropic exception are closely related to the three-dimensional
rigid-motion analysis of Tidor, Yu, and Zakharov
\cite[Sections~5--6]{tidoryuzakharov2026}.  Here the ruling geometry is
that of evaluation flats of alternating forms in arbitrary fixed
dimension; the required degree bound and exceptional-family
classification are proved in full below.

\subsection{Approximate complete intersections}
\label{aci:section}

We prove Walsh's approximate complete-intersection theorem
\cite[Theorem~1.3]{walsh2020}. The partial-degree framework and the
construction of compatible cuts follow
\cite[Definition~2.2 and Lemma~5.3]{walsh2020}.
The comparison of restriction spaces appears in the proof of
\cite[Lemma~3.5]{walsh2020}. Walsh obtains a corresponding prefix
degree estimate for his chosen cuts
\cite[Theorem~5.5]{walsh2020}.
We prove the forms needed here; the component Hilbert bounds in
Lemma~\ref{poly:hilbert} are the only quantitative algebraic input.
In particular, this proof is independent
of Theorem~\ref{conc:theorem}.

\begin{theorem}[Approximate complete intersection]
\label{aci:theorem}
\label{conc:walsh}
For every $n$ there is a constant $C_n$ with the following property.
If $V\subset\mathbb A^n_{\mathbb C}$ is irreducible of codimension
$c>0$, there are nonconstant polynomials $f_1,\ldots,f_c\in I(V)$ such
that $V$ is an irreducible component of $Z(f_1,\ldots,f_c)$ and
\[
 \prod_{i=1}^c\deg f_i\le C_n\deg V.
\]
\end{theorem}

Here and below, a component \emph{containing $V$} means that it contains
the whole irreducible variety $V$.  Set
\[
 I_{\le t}(V)=I(V)\cap\mathbb C[x_1,\ldots,x_n]_{\le t},
 \qquad
 H_V(t)=\dim_{\mathbb C}
 \bigl(\mathbb C[x_1,\ldots,x_n]_{\le t}/I_{\le t}(V)\bigr).
\]
For $1\le i\le c$, define $\delta_i$ to be the smallest integer $t$
such that every irreducible component of $Z(I_{\le t}(V))$ containing
$V$ has codimension at least $i$.  The notation $Z(I_{\le t}(V))$
means the zero set of a basis of this finite-dimensional vector space.
Finite generation of $I(V)$ shows that these integers exist, and their
definition gives
\begin{equation}
 1\le\delta_1\le\cdots\le\delta_c.
 \label{aci:partial-degrees}
\end{equation}
We also put $\delta_0=0$.

\begin{lemma}[Compatible partial cuts]
\label{aci:cuts}
There are $f_i\in I_{\le\delta_i}(V)$, $1\le i\le c$, such that
every component of $Z(f_1,\ldots,f_i)$ containing $V$ has codimension
exactly $i$.

For any such choice, if $0\le j<c$ and
$\delta_j\le m<\delta_{j+1}$, some component $W$ of
$Z(f_1,\ldots,f_j)$ containing $V$ satisfies
\begin{equation}
 I_{\le m}(W)=I_{\le m}(V),\qquad H_W(m)=H_V(m).
 \label{aci:shared-rank}
\end{equation}
When $j=0$, the prefix zero set and its unique component are
$\mathbb A^n$.
\end{lemma}

\begin{proof}
Suppose $f_1,\ldots,f_{i-1}$ have been chosen, and list the finitely
many components $W$ of their zero set which contain $V$.  Each has
codimension $i-1$.  None is contained in $Z(I_{\le\delta_i}(V))$:
otherwise a component of that latter set would contain $W$, hence
contain $V$ and have codimension at most $i-1$, contrary to the
definition of $\delta_i$.  Consequently the members of
$I_{\le\delta_i}(V)$ that vanish on any fixed $W$ form a proper
linear subspace.  Choose $f_i$ outside their finite union.  It cuts
every such $W$ properly.  The principal ideal theorem then gives
codimension exactly $i$ for every new component containing $V$.
This includes $i=1$, when the preceding zero set is the ambient space,
and proves the construction.

For the second assertion, suppose no component $W$ in the indicated
prefix has $I_{\le m}(V)\subset I(W)$.  The same finite-union argument
provides $g\in I_{\le m}(V)$ which vanishes identically on none of
them.  All components of $Z(f_1,\ldots,f_j,g)$ containing $V$ then
have codimension $j+1$.  Since $m\ge\delta_j$, every one of these
equations belongs to $I_{\le m}(V)$.  Thus every component of
$Z(I_{\le m}(V))$ containing $V$ has codimension at least $j+1$.
This contradicts $m<\delta_{j+1}$.  For a component $W$ where the
inclusion does hold, the reverse inclusion follows from $V\subset W$,
giving \eqref{aci:shared-rank}.
\end{proof}

\begin{proof}[Proof of Theorem~\ref{aci:theorem}]
Fix the polynomials from Lemma~\ref{aci:cuts}.  We will prove, for
$0\le s\le c$, that every component $U$ of the first $s$ cuts
containing $V$ satisfies
\begin{equation}
 \deg U\ge\gamma_n^{\,s}\prod_{i=1}^s\delta_i
 \label{aci:prefix-degree}
\end{equation}
for one sufficiently small $0<\gamma_n\le1$ depending only on $n$.
For $s=0$ both sides are $1$.

We record the precise consequence of Lemma~\ref{poly:hilbert} used
here.  There are constants $a_n>0$ and $b_n\ge1$ such that an
irreducible $k$-dimensional variety $X$ satisfies
\begin{equation}
 H_X(m)\le b_n\deg(X)m^k\qquad(m\ge1),
 \label{aci:rank-upper}
\end{equation}
and, if $X$ is a component of a system whose equations have degrees at
most $m$, then
\begin{equation}
 H_X(m)\ge a_n\deg(X)m^k\qquad(m\ge1).
 \label{aci:rank-lower}
\end{equation}
The second assertion is the bound at the equation degree in
\eqref{poly:hilbert-at-L}, applied with $L=m$.  It permits other
components of the system to have larger dimension.  The lower bound
also holds for $X=\mathbb A^n$, after decreasing $a_n$, directly from
$H_{\mathbb A^n}(m)=\binom{m+n}{n}$.

Suppose \eqref{aci:prefix-degree} is known for smaller indices, and
let $U$ be a component at index $s$.  If $\delta_s=1$, the entire
product in \eqref{aci:prefix-degree} is $1$, and $\deg U\ge1$ proves
the claim.  Otherwise set
\[
 m=\delta_s-1,\qquad
 j=\max\{i:0\le i<s,\ \delta_i<\delta_s\}.
\]
Thus
\begin{equation}
 \delta_j\le m<\delta_{j+1}
   =\cdots=\delta_s\le2m.
 \label{aci:block}
\end{equation}
Lemma~\ref{aci:cuts} supplies a component $W$ at index $j$ with
$H_W(m)=H_V(m)$.  Its dimension is $n-j$, and the degrees of its
defining prefix equations are at most $\delta_j\le m$.
The induction hypothesis and \eqref{aci:rank-lower} therefore give
\[
 H_V(m)=H_W(m)
 \ge a_n\gamma_n^{\,j}
       \left(\prod_{i=1}^j\delta_i\right)m^{n-j}.
\]
This formula also applies when $j=0$.  Since $V\subset U$, restriction
from $U$ to $V$ is surjective on polynomial restriction spaces, so
$H_U(m)\ge H_V(m)$.  On the other hand, $\dim U=n-s$ and
\eqref{aci:rank-upper} bounds $H_U(m)$ above.  Combining the bounds,
then using \eqref{aci:block}, yields
\[
 \begin{split}
 \deg U
 &\ge \frac{a_n}{b_n}\gamma_n^{\,j}
       \left(\prod_{i=1}^j\delta_i\right)m^{s-j}\\
 &\ge \frac{a_n}{2^n b_n}\gamma_n^{\,j}
       \prod_{i=1}^s\delta_i.
 \end{split}
\]
Choose $\gamma_n\le\min\{1,a_n/(2^n b_n)\}$.  Since $j<s$, the
last expression is at least the right side of
\eqref{aci:prefix-degree}, completing the induction.  Notice that no
containment relation between $U$ and the chosen $W$ was needed; both
contain $V$.

At $s=c$, every prefix component containing $V$ has the same dimension
as $V$, so $V$ itself is such a component.  Applying
\eqref{aci:prefix-degree} to it gives
\[
 \prod_{i=1}^c\deg f_i
 \le\prod_{i=1}^c\delta_i
 \le\gamma_n^{-c}\deg V
 \le\gamma_n^{-n}\deg V.
\]
This proves the theorem with $C_n=\gamma_n^{-n}$.
\end{proof}

The conclusion concerns an irreducible component.  It does not assert
that the displayed equations generate its radical ideal or have
independent differentials; the next regularity argument supplies the
additional control needed here.


\subsection{Separators and regular equations}

We use the isolated-intersection form of B\'ezout, with local
multiplicities, and the standard local algebra and projective-bundle
facts detailed in the proofs below
\cite{fulton1998,hartshorne1977}.

\begin{lemma}[Counting by separators]
\label{conc:separators}
Under the separator assumption of Theorem~\ref{conc:theorem}, an
irreducible projective variety $Y\subset\mathbf P(E\oplus E^*)$ of
dimension $a$ contains at most
\[
 C_{e,C_0}\deg(Y)A^a
\]
sample parameters.  The variety $Y$ need not be contained in
$\mathcal Q$.
\end{lemma}

\begin{proof}
Induct on $a$.  The assertion for $a=0$ says that a point contains at
most one sample parameter.  Otherwise, if $Y\cap\mathcal P\ne\varnothing$,
choose $\lambda\in Y\cap\mathcal P$ and a separator $q$ for it.  Since
$q(\lambda)\ne0$, the section $Y\cap Z(q)$ is proper, and it contains
all the other counted parameters.  Its components have dimension
$a-1$ and total degree at most $C_0A\deg Y$.  Applying the induction
hypothesis to those components gives
\[
 |Y\cap\mathcal P|
 \leq 1+C_{a-1}C_0\deg(Y)A^a.
\]
The initial $1$ is absorbed because $A\geq1$ and $\deg Y\geq1$.
There are only boundedly many possible dimensions.
\end{proof}

\begin{lemma}[Regular equations of controlled degree]
\label{conc:regularity}
Suppose that an irreducible $V\subset\mathbb A^n_{\mathbb C}$ has
codimension $c>0$ and is a component of $Z(f_1,\ldots,f_c)$, where
\[
 \deg f_i\leq D,\qquad
 \prod_i\deg f_i\leq C_n\deg V,\qquad D\geq1,
\]
with the dimensional constant from Theorem~\ref{conc:walsh}.
There are polynomials $g_1,\ldots,g_c\in I(V)$ of degrees $O_n(D)$
whose differentials are independent at the generic point of $V$.
Moreover, a polynomial $H_0$ of degree $O_n(D)$, nonzero on $V$, can be
chosen so that $V\setminus Z(H_0)$ is smooth and has tangent space
\[
 T_zV=\bigcap_{i=1}^c\ker dg_i(z)
 \qquad(z\in V\setminus Z(H_0)).
\]
\end{lemma}

\begin{proof}
Put $P=I(V)$ and consider the regular local ring
\[
 R=\mathbb C[x_1,\ldots,x_n]_P,\qquad
 \mathfrak m=PR,\qquad K=R/\mathfrak m=\mathbb C(V).
\]
Its dimension is $c$.  Because $V$ is an isolated component of the
common zero set at its generic point, $J=(f_1,\ldots,f_c)R$ is
$\mathfrak m$-primary.  Its generators form a parameter sequence,
which is regular because $R$ is Cohen--Macaulay
\cite[Tags~00OQ, 00NQ, and 02JN]{stacksproject}.
We first bound the length of $R/J$ by a dimensional constant.

Write $k=\dim V$.  A general projective linear projection of
$\overline V$ to $\mathbf P^k$, with center in the hyperplane at
infinity disjoint from $\overline V$, is finite and has degree
$\deg V$.  Finiteness follows because a positive-dimensional projective
fiber would meet its hyperplane lying in the projection center.
Restricting to the affine chart and changing affine coordinates gives
base coordinates $t_1,\ldots,t_k$ and complementary coordinates
$w_1,\ldots,w_c$.  Over $F=\mathbb C(t_1,\ldots,t_k)$, the generic
fiber of $V$ is a closed point of $\mathbb A^c_F$ with residue field
$K$, where $[K:F]=\deg V$.  Localizing $F[w_1,\ldots,w_c]$ at that
point gives the same ring $R$.

After extending $F$ to an algebraic closure, this point splits into
$[K:F]$ points, since the extension is separable.  The sum of their
local intersection multiplicities for the equations $f_i$ is
$[K:F]\operatorname{length}_R(R/J)$.  They are isolated common zeros;
other components of the system do not affect their isolated B\'ezout
bound \cite[Theorem~12.3 and Example~12.3.1]{fulton1998}.
The equations $f_i$ define hypersurfaces in regular local ambient
rings, so these factors are Cohen--Macaulay and the isolated
complete-intersection multiplicity equals the quotient length;
see \cite[Example~12.3.7(ii)--(iii)]{fulton1998} for this identification
and its Cohen--Macaulay hypothesis.  Hence
\begin{equation}
 \operatorname{length}_R(R/J)
 \leq\frac{\prod_i\deg f_i}{\deg V}\leq C_n.
 \label{conc:local-length}
\end{equation}
Choose an integer $L=O_n(1)$ at least this length.  The maximal ideal
of an Artinian local ring of length at most $L$ has $L$th power zero,
so
\begin{equation}
 \mathfrak m^L\subset J.
 \label{conc:adic-bound}
\end{equation}

We construct the $g_i$ one at a time.  Suppose $g_1,\ldots,g_u$ have
independent classes in $\mathfrak m/\mathfrak m^2$, where $u<c$.
They are part of a regular parameter system, and therefore
\[
 R_u=R/(g_1,\ldots,g_u),\qquad
 \mathfrak m_u=\mathfrak mR_u
\]
is a regular local ring of positive dimension.  The images of the
$f_i$ generate an ideal containing $\mathfrak m_u^L$.  At least one
image has finite $\mathfrak m_u$-adic order $b$ with $1\leq b\leq L$.
Otherwise all the images would lie in $\mathfrak m_u^{L+1}$, whence
$\mathfrak m_u^L=\mathfrak m_u^{L+1}$ and Nakayama's lemma would give
$\mathfrak m_u^L=0$, impossible in this positive-dimensional regular
local ring.

We need derivatives which do not disturb the already chosen equations.
For $u=0$ use coordinate derivations.  For $u>0$, choose a $u$-square
Jacobian block $J_0$ of $(g_1,\ldots,g_u)$ whose determinant $\Delta$
is nonzero in $K$.  If $J_t$ is any remaining column, form a polynomial
vector field having coefficient $\Delta$ in coordinate $t$ and
coefficients $-\operatorname{adj}(J_0)J_t$ in the chosen coordinates.
These $n-u$ fields, denoted by $\delta_\alpha$, satisfy
\[
 \delta_\alpha g_i=0\quad(1\leq i\leq u)
\]
as polynomial identities.  Thus they extend to $R$ and descend to
$R_u$.  If the $g_i$ have maximum degree $E_u$, the coefficient degrees
of these fields are $O_n(E_u)$.

At the generic point, the fields span the common kernel of the
$d g_i$.  The conormal map
$\mathfrak m/\mathfrak m^2\to K^n$, $f\mapsto df$, is injective there,
as $V$ is generically smooth in characteristic zero.  It follows that
the residue evaluations of the $\delta_\alpha$ span the dual of
$\mathfrak m_u/\mathfrak m_u^2$.  These residue evaluations are
$K$-linear: the extra Leibniz term from differentiating a coefficient
has a factor in $\mathfrak m_u$.

Take an $f_i$ of order $b$ as above.  A derivation lowers adic order by
at most one, so fewer than $b$ of these derivations still give zero
residue.  Some composition of $b$ of them has nonzero residue.  Indeed,
the leading term of $f_i$ belongs to the degree-$b$ piece of the graded
ring of $R_u$ \cite[Tag~00NO]{stacksproject}, namely
\[
 \mathfrak m_u^b/\mathfrak m_u^{b+1}
 \simeq\operatorname{Sym}^b_K(\mathfrak m_u/\mathfrak m_u^2).
\]
Its polarization is nonzero in characteristic zero.  Since the residue
functionals span the dual, some $b$-tuple detects it.  In computing
that residue, terms in which a derivation hits a coefficient retain a
positive-order factor and vanish.  This gives exactly the asserted
iterated derivative.

Stop immediately before the last differentiation and call the resulting
polynomial $g_{u+1}$.  Its residue in $R_u$ is zero, so it belongs to
$P$.  Its derivative has nonzero residue, so its class in
$\mathfrak m_u/\mathfrak m_u^2$ is nonzero.  It adds an independent
conormal.  This procedure uses at most $L-1$ derivatives per stage and
at most $c$ stages.  Applying a polynomial field of coefficient degree
$T$ to a degree-$B$ polynomial gives degree at most $B+T-1$; the bounds
on $L$ and $c$ therefore keep every $g_i$ of degree $O_n(D)$.

Finally choose a $c$-square minor $H_0$ of the Jacobian of the $g_i$
which is nonzero on $V$.  Its degree is $O_n(D)$.  Where it is nonzero,
the common zero set of the $g_i$ is locally smooth of dimension $n-c$.
The contained variety $V$ has the same dimension, so its germ equals
that smooth germ.  This proves the assertion about tangent spaces.
\end{proof}

\subsection{Recovering actual rulings from tangent directions}

In a chart $v_a=1$, a \emph{ruling incidence} is an irreducible set of
pairs $(z,v)$ for which $z+K_v\subset V$.  Its parameter map is
\begin{equation}
 (z,v)\longmapsto[v:zv].
 \label{conc:parameter-map}
\end{equation}
The point of the next lemma is that both the number of such incidences
and the degree of their exceptional set can be controlled.

\begin{lemma}[Controlled ruling incidences]
\label{conc:rulings}
Let $V\subset S$ be irreducible of dimension $k=s_0+h$, with
$1\leq h<r$.  Suppose that $V$ has $c=n-k$ approximate equations as
in Lemma~\ref{conc:regularity}, of maximum degree at most $D$.
There is a polynomial $H$, nonzero on $V$ and of degree $O_e(D)$,
such that, on $\Omega=V\setminus Z(H)$, the following hold.
\begin{enumerate}
\item There are $O_e(1)$ smooth irreducible ruling incidences $\Gamma$.
      Each is closed in $\Omega\times\{v_a=1\}$ for some $a$ and maps
      smoothly and dominantly to $\Omega$.
\item The fibers over $\Omega$ are either finite of bounded size or
      affine lines of normalized vectors $v$.  Denote the fiber
      dimension by $j=0$ or $j=1$.  The latter is possible only if
      $c=1$.
\item Every flat $F_\lambda\subset V$ meeting $\Omega$ is represented
      by at least one incidence.  All other contained flats lie in
      $V\cap Z(H)$.
\item If $\Lambda$ is the projective closure of the parameter image
      of one incidence, then $\dim\Lambda=h+j$.  Every nonempty
      parameter fiber is exactly $F_\lambda\cap\Omega$ and has
      dimension $s_0$.
\end{enumerate}
\end{lemma}

\begin{proof}
Use Lemma~\ref{conc:regularity} to choose $g_1,\ldots,g_c$ and an
initial smooth open set in $V$.  In the construction below we remove
zeros of a bounded number of polynomials, each of degree $O_e(D)$ and
nonzero on $V$.  Their product, together with the initial Jacobian
minor, will be $H$.  We continue to call the current open set $\Omega$.

First consider tangent directions.  Under $S^*=\bigwedge^2E$, the
annihilator of $K_v$ is $v\wedge E$.  A bivector belongs to this space
if and only if its wedge with $v$ is zero.  Thus the necessary
condition $K_v\subset T_zV$ is the linear system
\begin{equation}
 v\wedge dg_i(z)=0\qquad(1\leq i\leq c).
 \label{conc:tangent-equations}
\end{equation}
Its vector solution space has dimension at most two.  To see this,
let a nonzero bivector $w$ satisfy $v\wedge w=0$ for a nonzero $v$.
Splitting off $\mathbb Cv$ shows that $w=v\wedge a$ for some $a$;
the vectors $v'$ satisfying $v'\wedge w=0$ are precisely
$\operatorname{span}(v,a)$.  If two independent conormals had a common
two-dimensional solution space, both would lie in its one-dimensional
exterior square, a contradiction.  A two-dimensional solution space
therefore forces $c=1$.

Work separately in each of the $e$ normalized charts $v_a=1$.
Gaussian elimination over $\mathbb C(V)$, followed by removal of
nonzero rank minors, puts the system throughout $\Omega$ into one of
three forms: no solution, a graph $v=b(z)$, or
\begin{equation}
 v=b(z)+\tau a_1(z).
 \label{conc:direction-graph}
\end{equation}
In the last case choose a free coordinate of $v$ as $\tau$, so that
its coordinate in $b$ is zero and in $a_1$ is one.  All coefficients
are regular on $\Omega$ and admit ambient rational representatives
with numerator and denominator degrees $O_e(D)$, since the matrices
have bounded size.  Rank conditions which vanish generically vanish
on all of $V$; excluding the nonzero minors consequently controls all
remaining fibers.  Denote the resulting solution variety by $\Gamma_0$.
It is a smooth graph over $\Omega$, or a smooth graph over
$\Omega\times\mathbb A^1$.

We now test whether the tangent directions integrate to a contained
flat.  Write $\varepsilon_1,\ldots,\varepsilon_e$ for a basis of $E^*$.
In the chart $v_a=1$, exterior products of the covectors
$\varepsilon_i-v_i\varepsilon_a$, $i\ne a$, form a basis
$B_1(v),\ldots,B_{s_0}(v)$ of $K_v$.  Their coefficients have degree
at most two in $v$.  Consider the horizontal fields
\[
 X_\alpha(z,v)=(B_\alpha(v),0).
\]
Their $z$-components are tangent to $V$, by
\eqref{conc:tangent-equations}.  They are tangent to $\Gamma_0$ if and
only if the derivatives of all coordinates of its graph formula
vanish in direction $B_\alpha(v)$, with $\tau$ fixed in
\eqref{conc:direction-graph}.  Holding $\tau$ fixed is required
because it is one of the coordinates of the normalized vector $v$.
Call these derivative expressions $T_\beta$.

With no free coordinate the $T_\beta$ are rational functions of $z$.
With a free coordinate they are polynomials in $\tau$ of bounded
degree, whose coefficients have ambient rational degree $O_e(D)$.
This follows by substituting \eqref{conc:direction-graph} in the
quadratic expressions $B_\alpha$ and differentiating its rational
coefficients.  Differentiation along $T_zV$ is independent of the
ambient representatives chosen, since it annihilates functions
vanishing on $V$.

Every contained flat passes this first test: on a small open part of
the flat, $v$ is fixed and the direction equations continue to hold.
If every $T_\beta$ is identically zero, retain $\Gamma_0$; its
horizontal fields are already tangent everywhere.  If there is no
free coordinate but one test is nonzero, remove the zero set of its
numerator and discard this chart.  No contained flat in that chart
can meet the remaining open set.

It remains to treat a free coordinate and a nonidentically zero test.
Let $P(\tau)$ be the monic gcd of the tests over $\mathbb C(V)$, and
let $G(\tau)$ be its monic squarefree part.  The Euclidean algorithm
gives identities
\[
 T_\beta=P Q_\beta,\qquad
 P=\sum_\beta U_\beta T_\beta,
 \qquad G\mid P,\quad P\mid G^q
\]
for a bounded integer $q$.  Exclude the poles of all coefficients
used in these identities.  They then show in every remaining fiber
that the common roots of the tests are exactly the roots of $G$.
If $G=1$, the chart contributes nothing.  Otherwise exclude also the
zero set of its discriminant.  The root variety
\[
 \Sigma_G=\{(z,\tau)\in\Omega\times\mathbb A^1:G(z,\tau)=0\}
\]
is now finite \emph{\'etale} over $\Omega$, with bounded degree.  In
particular it is smooth, and its roots form local holomorphic sheets.

All these exclusions have total degree $O_e(D)$.  Indeed, the number
of tests and their degrees in $\tau$ are bounded.  Gcd, squarefree
part, quotient, and discriminant computations therefore use a bounded
number of field operations.  A bounded number of operations on
ambient fractions of degree $O_e(D)$ preserves that degree bound.
Every excluded denominator or pivot has nonzero restriction to $V$.
The same observation applies to the following second test.

Along a contained flat, $G(z,\tau)=0$ persists with $\tau$ fixed.
Thus a second necessary condition is
\[
 R_\alpha(z,\tau)
 =\sum_{m=1}^{n}(B_\alpha(v(z,\tau)))_m
                  \frac{\partial G}{\partial z_m}(z,\tau)=0.
\]
These expressions again have bounded degree in $\tau$ and coefficient
degrees $O_e(D)$.  Form the monic polynomial
\[
 P_2=\gcd(G,R_1,\ldots,R_{s_0})
\]
and exclude the poles needed to specialize its quotient and B\'ezout
identities.  Its root variety $\Sigma_{P_2}$ is empty or finite
\emph{\'etale} over $\Omega$.  Because $G$ remains squarefree and
$P_2$ divides it, the factorization $G=P_2(G/P_2)$ has disjoint root
sets in every remaining fiber.  Thus $\Sigma_{P_2}$ is locally a union
of whole sheets of $\Sigma_G$.

At a point of $\Sigma_{P_2}$, the first test makes $X_\alpha$ tangent
to $\Gamma_0$, where its coordinates are $(B_\alpha(v),0)$ in
$(z,\tau)$.  The second test makes it tangent to $G=0$.  Hence it is
tangent to the selected sheet of $\Sigma_{P_2}$.  No further
sequence of tangency tests is needed.  Retain its irreducible
components as the incidences.  A smooth variety has disjoint
irreducible components.  Each component of this finite \emph{\'etale}
cover maps openly and closedly to the irreducible base, hence
surjectively; their number and fiber sizes are bounded by $\deg G$.

We verify containment rather than merely tangency.  On every retained
incidence the fields $X_\alpha$ are tangent.  Their ambient flows are
\[
 (z,v)\longmapsto(z+tB_\alpha(v),v),
\]
since $v$ is constant.  Uniqueness of local holomorphic flows keeps
these translations in the incidence for small $t$.  Composing them
produces an analytic open subset of $z+K_v$, because the $B_\alpha(v)$
form a basis of $K_v$.  Every polynomial vanishing on $V$ vanishes on
that open subset, hence on the whole affine flat.  Thus every retained
point really represents a contained flat.

After processing all normalized charts, intersect their open sets and
restrict their incidences to this common $\Omega$.  The product of
the excluded polynomials is nonzero on irreducible $V$ and has degree
$O_e(D)$.  Every nonempty restricted incidence remains irreducible and
dominant.  Its fibers are finite of bounded size, except for the
retained affine-line graphs; the latter arose only in codimension one.

If a contained flat meets $\Omega$, choose a chart for its $v$ and a
point of $F_\lambda\cap\Omega$.  Both tests are necessary along that
flat, and the specialized gcd identities therefore place the point in
a retained incidence.  This proves coverage.

Finally fix one incidence and a nonempty fiber of
\eqref{conc:parameter-map}.  Its normalized vector $v$ is fixed, and
the fiber is closed in $F_\lambda\cap\Omega$.  The preceding flows
preserve both $v$ and $zv$ and give an analytic open subset of that
flat in the fiber.  Such an open subset is Zariski dense in the
irreducible variety $F_\lambda\cap\Omega$, so the fiber equals this
whole intersection.  It has dimension $s_0$.  The fiber-dimension
theorem now gives
\[
 \dim\Lambda=\dim\Gamma-s_0=(k+j)-s_0=h+j.
\]
\end{proof}

\subsection{Degree bounds and exceptional parameter families}
\label{conc:deg-section}

The ruling construction bounds the number of incidences and the degree
of the set removed from \(V\). For an incidence with \(j=0\), its
projective parameter closure \(\Lambda\) has dimension \(h\), so
Lemma~\ref{conc:separators} counts at
most \(C\deg(\Lambda)A^h\) sample parameters. A bound
\(\deg\Lambda\le C_e\deg V\) would therefore give the scale required
by Theorem~\ref{conc:theorem}. When \(j=1\), the parameter variety has
dimension \(h+1\), and the separator bound alone would cost one extra
factor of \(A\). The next proposition proves that this case has
\(h=r-1\) and that the parameter variety is a maximal isotropic
\(\mathbf P^r\subset\mathcal Q\), where assumption (ii) gives the
required \(O(A^{r-1})\) bound. The projective incidence bundle
constructed below gives the degree estimate and the tangent constraint
used to classify this exceptional case.

\begin{proposition}
\label{conc:degree}
Let \(V\subsetneq S\) be irreducible of dimension \(k=s_0+h\), where
\(1\le h<r\), and let \(\Gamma\) be a ruling incidence provided by
Lemma~\ref{conc:rulings}.  Write \(j\in\{0,1\}\) for the dimension
of its fibers over the working open subset \(\Omega\subset V\).
Let \(\Lambda\) be the projective closure of its parameter image
\[
 (z,v)\longmapsto[v:zv]
 \quad\hbox{in}\quad
 \mathcal Q=\{[v:y]\in\mathbf{P}(E\oplus E^*):y(v)=0\}.
\]
Then
\[
 \dim\Lambda=h+j,
 \qquad \deg\Lambda\le C_e\deg V.
\]
If \(j=1\), then \(h=r-1\) and \(\Lambda\) is a linear
\(\mathbf{P}^r\subset\mathcal Q\), hence a maximal isotropic
projective subspace.
\end{proposition}

\begin{proof}
The fibers of the parameter map are the nonempty sets
\(F_\lambda\cap\Omega\), of dimension \(s_0\), whereas
\(\dim\Gamma=k+j\).  The fiber dimension theorem therefore gives
\(\dim\Lambda=h+j\).  We prove the degree estimate and the
assertion for \(j=1\) in turn.

\paragraph{Compactifying the parameter space.}
On \(\mathbf{P}(E)\), let \(L=\mathcal O_{\mathbf{P}(E)}(-1)\)
be the tautological line bundle and put
\(Q=(E\otimes\mathcal O)/L\).  Evaluation of a skew form on
the line \(L\) gives the exact sequence
\begin{equation}
 0\longrightarrow K\longrightarrow S\otimes\mathcal O
 \longrightarrow Q^*\otimes L^*\longrightarrow0,
 \qquad K=\bigwedge^2 Q^*.
 \label{conc:deg-evaluation}
\end{equation}
Indeed, the value of a skew form on \(L\) is a map from \(L\)
to its annihilator \(L^\perp\subset E^*\), every such map
extends to a skew form on \(E\), and the kernel consists of the
forms on \(E/L\).  The ranks of these three bundles are
\(s_0,n,r\), respectively.

All projective bundles below parametrize lines.  Consider
\begin{equation}
 B=\mathbf{P}\bigl(\mathcal O\oplus Q^*\otimes L^*\bigr)
   \simeq\mathbf{P}(L\oplus Q^*),
 \qquad \rho:B\longrightarrow\mathbf{P}(E).
 \label{conc:deg-parameter-bundle}
\end{equation}
The second presentation is obtained by tensoring the first bundle
with \(L\).  We use the line convention of
\cite[Section~3.1 and Example~3.1.1]{fulton1998}.  Since \(L\oplus Q^*\)
is a subbundle of the trivial bundle with fiber \(E\oplus E^*\),
it induces a morphism \(\mu:B\to\mathcal Q\).  Explicitly,
write a point in the first presentation as \((L,[t:\psi])\),
with \(\psi:L\to L^\perp\).  For any nonzero \(v\in L\),
\begin{equation}
 \mu(L,[t:\psi])=[tv:\psi(v)].
 \label{conc:deg-parameter-map}
\end{equation}
Both rescaling \((t,\psi)\) and replacing \(v\) by another
generator of \(L\) leave this point unchanged.  Its quadratic
value is zero because \(\psi(v)\) annihilates \(v\).

Over \(\mathcal Q^\circ=\{[v:y]\in\mathcal Q:v\ne0\}\),
the inverse to \(\mu\) recovers \(L=\mathbb{C}v\) and
\([1:\psi]\), where \(\psi(v)=y\).  Thus \(\mu\) is an
isomorphism over this open set.  A boundary point \([0:y]\)
has fiber \(\mathbf{P}(\ker y)\): its preimages have \(t=0\),
and \(L\) can be any line annihilated by \(y\).

Let \(\mathcal T=\mathcal O_B(-1)\) be the tautological line
in the \emph{first} presentation of
\eqref{conc:deg-parameter-bundle}, and set
\[
 a=\rho^*c_1\bigl(\mathcal O_{\mathbf{P}(E)}(1)\bigr),
 \qquad b=c_1(\mathcal T^*),
 \qquad \ell=\mu^*c_1\bigl(\mathcal O_{\mathcal Q}(1)\bigr).
\]
The tautological line in the second presentation is
\(\mathcal T\otimes\rho^*L\).  Consequently
\begin{equation}
 \ell=a+b.
 \label{conc:deg-hyperplane}
\end{equation}
The classes \(a\) and \(\ell\) are globally generated.  We do
not need positivity of \(b\).

\paragraph{The projective incidence and its boundary.}
Pull back \eqref{conc:deg-evaluation} to \(B\).  Define
\(\mathcal U\) as the inverse image of \(\mathcal T\)
under the surjective bundle map
\[
 (\mathbb{C}\oplus S)\otimes\mathcal O_B
 \longrightarrow
 \mathcal O_B\oplus\rho^*(Q^*\otimes L^*),
 \qquad (x_0,z)\longmapsto(x_0,z|_L).
\]
It is a vector bundle of rank \(s_0+1\), with exact sequence
\begin{equation}
 0\longrightarrow\rho^*K\longrightarrow\mathcal U
 \longrightarrow\mathcal T\longrightarrow0.
 \label{conc:deg-incidence-bundle}
\end{equation}
Its inclusion in the trivial bundle defines a morphism
\(\mathbf{P}(\mathcal U)\to\mathbf{P}(\mathbb{C}\oplus S)\).
The pullback of the hyperplane class under this morphism is
\[
 \xi=c_1\bigl(\mathcal O_{\mathbf{P}(\mathcal U)}(1)\bigr).
\]
For a parameter represented by \((L,[1:\psi])\), choose
\(v\in L\setminus\{0\}\) and write \(y=\psi(v)\).
The projective fiber maps isomorphically to
\begin{equation}
 \{[x_0:z]:zv=x_0y\}=\overline{F_{[v:y]}}.
 \label{conc:deg-flat-closure}
\end{equation}
Over \((L,[t:\psi])\), the defining condition is instead
\((x_0,z|_L)\in\mathbb{C}(t,\psi)\).  In particular,
\begin{equation}
 t=0\quad\Longrightarrow\quad x_0=0.
 \label{conc:deg-boundary}
\end{equation}
Thus no boundary parameter occurs over an affine point \([1:z]\).

Lift the parameter image of \(\Gamma\) to \(B\) using the
isomorphism over \(\mathcal Q^\circ\), and denote its reduced
irreducible closure by \(\widetilde\Lambda\).  The map
\(\mu:\widetilde\Lambda\to\Lambda\) is birational.  Put
\[
 X=\mathbf{P}(\mathcal U|_{\widetilde\Lambda}),
 \qquad \pi:X\longrightarrow\widetilde\Lambda.
\]
The variety \(X\) is irreducible and has dimension
\(s_0+h+j=k+j\).  The point coordinate and the lifted parameter
give an embedding of \(\Gamma\) as a locally closed subset of
\(X\): in its normalized direction chart the inverse recovers
the original pair \((z,v)\).  Since the dimensions agree, this
subset is dense and contains a dense open subset of \(X\).
The projective point map on \(X\) therefore has image exactly
\(\overline V\).  More explicitly, its image is closed by
properness; density shows that it is contained in the closure of
the point image of \(\Gamma\), and dominance of
\(\Gamma\to\Omega\) gives the opposite inclusion.  We obtain
a surjective morphism
\[
 f:X\longrightarrow\overline V,
 \qquad
 \xi=f^*c_1\bigl(\mathcal O_{\overline V}(1)\bigr).
\]
Here and below we also use \(a,b,\ell\) for their restrictions
and pullbacks to the varieties on which they occur.

\paragraph{The degree of a generic point fiber.}
Fix an affine \(z\in S\).  By \eqref{conc:deg-boundary}, the
full incidence over \([1:z]\), before restriction to
\(\widetilde\Lambda\), has parameters \([v:zv]\) with
\(v\ne0\), and is isomorphic to \(\mathbf{P}(E)\).
On this projective space \(\ell\) is the ordinary hyperplane
class: the map \([v]\mapsto[v:zv]\) is induced by an injective
linear map \(E\to E\oplus E^*\).

Choose a dense open subset \(X^\circ\subset X\) lying in the
lift of \(\Gamma\).  Every component of its complement has
dimension at most \(k+j-1\).  A component which dominates
\(V\) therefore has generic fiber dimension at most \(j-1\);
a component which does not dominate \(V\) is absent over a
general point.  The complement cannot supply a top-dimensional
component of the generic fiber of \(f\).

When \(j=0\), the generic fiber thus consists of at most \(C_e\)
points, as does the fiber of \(\Gamma\to\Omega\).  When
\(j=1\), the generic fiber of \(\Gamma\to\Omega\) is an
affine line of normalized vectors.  Its closure in \(\mathbf{P}(E)\)
is a projective line, of degree one against \(\ell\), and this
is the only top-dimensional generic fiber component of \(f\).
These components have multiplicity one.  Indeed, the fibers on
the dense open ruling incidence are smooth: its projection is
finite \'etale or a smooth affine-line family.  This determines
the generic multiplicity of each component, which cannot be
changed by a complement of smaller fiber dimension.

If \(d_\Gamma\) denotes the degree of the geometric generic
fiber against \(\ell^j\), we have proved
\begin{equation}
 d_\Gamma\le C_e,
 \qquad d_\Gamma=1\quad\hbox{when }j=1.
 \label{conc:deg-fiber-degree}
\end{equation}
Proper pushforward and the projection formula, applied to
fundamental cycles, now give
\begin{equation}
 f_*(\ell^j\cap[X])=d_\Gamma[\overline V],
 \qquad
 \int_X\ell^j\xi^k=d_\Gamma\deg V\le C_e\deg V.
 \label{conc:deg-point-integral}
\end{equation}
The coefficient in the first identity is computed on the generic
fiber.  Any other component of a representative hyperplane
intersection has image of dimension less than \(k\) and pushes
forward to zero.  These formulas do not require smoothness of
\(X\) or \(\overline V\); see
\cite[Sections 1.4, 2.3, 2.5]{fulton1998}.

\paragraph{A positive Segre-class formula.}
For projective bundles of lines, the projective bundle formula is
\begin{equation}
 \pi_*(\xi^{s_0+h}\cap[X])
   =s_h(\mathcal U|_{\widetilde\Lambda})
      \cap[\widetilde\Lambda],
 \qquad s_t(\mathcal U)=c_t(\mathcal U)^{-1}.
 \label{conc:deg-projective-bundle}
\end{equation}
This formula applies equally to a vector bundle on a singular
base \cite[Sections 3.1--3.3]{fulton1998}.  We compute it before
restriction to \(\widetilde\Lambda\).

Dualizing the tautological sequence on \(\mathbf{P}(E)\) and
tensoring by \(L^*=\mathcal O(1)\) gives
\[
 0\longrightarrow Q^*(1)
 \longrightarrow E^*\otimes\mathcal O(1)
 \longrightarrow\mathcal O(2)\longrightarrow0.
\]
Together with \eqref{conc:deg-evaluation}, this implies
\[
 c_t(K)^{-1}=c_t(Q^*(1))=\frac{(1+at)^e}{1+2at}.
\]
Since \(c_t(\mathcal T)=1-bt\), Equations
\eqref{conc:deg-incidence-bundle} and
\eqref{conc:deg-hyperplane} yield
\begin{equation}
 s_t(\mathcal U)
   =\frac{(1+at)^e}{(1+2at)(1-bt)}
   =\frac{(1+at)^e}{(1+2at)(1+(a-\ell)t)}.
 \label{conc:deg-segre-series}
\end{equation}
Expanding
\[
 \frac{1}{1+(a-\ell)t}
   =\sum_{i\ge0}\frac{\ell^i t^i}{(1+at)^{i+1}}
\]
and taking the coefficient of \(t^h\) gives
\begin{equation}
 s_h(\mathcal U)
 =\sum_{i=0}^h \ell^i a^{h-i}c(e-1-i,h-i),
 \qquad
 c(p,u)=[z^u]\frac{(1+z)^p}{1+2z}.
 \label{conc:deg-segre-coefficients}
\end{equation}

For integers \(p\ge u+1\) and \(u\ge0\), these scalar
coefficients are nonnegative.  Indeed,
\(c(p,0)=1\), and \(c(p,1)=p-2\ge0\) in its asserted
range.  For \(u\ge2\), the identity
\((1+z)^2=(1+2z)+z^2\) gives
\[
 c(p,u)=\binom{p-2}{u}+c(p-2,u-2).
\]
The second term lies in the same range, because
\(p-2\ge(u-2)+1\).  Induction proves the claim, with a
binomial coefficient equal to zero when its lower index exceeds
its upper index.  In \eqref{conc:deg-segre-coefficients}, the
required inequality is
\(e-1-i\ge h-i+1\), which follows from
\(h\le r-1=e-2\).  The coefficient of \(\ell^h\) is
\(c(e-1-h,0)=1\).

All mixed top intersections of \(a\) and \(\ell\) on
\(\widetilde\Lambda\) are nonnegative.  They can be computed
by successive general members of globally generated linear
systems, whose intersections with the fundamental cycle are
effective or empty.  Birationality of
\(\widetilde\Lambda\to\Lambda\), followed by
\eqref{conc:deg-segre-coefficients},
\eqref{conc:deg-projective-bundle}, and
\eqref{conc:deg-point-integral}, therefore gives
\begin{equation}
 \deg\Lambda
 =\int_{\widetilde\Lambda}\ell^{h+j}
 \le\int_{\widetilde\Lambda}\ell^j s_h(\mathcal U)
 =\int_X\ell^j\xi^k
 \le C_e\deg V.
 \label{conc:deg-bound}
\end{equation}

\paragraph{The tangent constraint in the affine-line case.}
Suppose now that \(j=1\).  Lemma~\ref{conc:rulings} gives
\(\operatorname{codim}_S V=1\), so \(h=r-1\) and
\(\dim\Lambda=r\).  We first show that the tangent space of
\(\Lambda\) at a general smooth point is maximal isotropic for
the tangent quadratic form of \(\mathcal Q\).

Choose a general smooth parameter \(\lambda=[v:y]\in
\Lambda\cap\mathcal Q^\circ\), fix \(z_0\in F_\lambda\),
and let \(U=E/\mathbb{C}v\), of dimension \(r\).  With a
representative \((v,y)\) fixed, tangent vectors to \(\mathcal Q\)
are pairs \((\dot v,\dot y)\) satisfying
\(\dot y(v)+y(\dot v)=0\), modulo \(\mathbb{C}(v,y)\).
There is an isomorphism
\begin{equation}
 T_\lambda\mathcal Q\longrightarrow U\oplus U^*,
 \qquad
 (\dot v,\dot y)\longmapsto
 (\alpha,\beta)=(\dot v\bmod v,\dot y-z_0\dot v).
 \label{conc:deg-tangent-coordinates}
\end{equation}
The covector \(\beta\) annihilates \(v\), by skew symmetry
and \(z_0v=y\).  Adding a multiple of \((v,y)\) changes
neither coordinate, and the kernel before quotienting is exactly
that line, proving the isomorphism.  The tangent quadratic form,
defined up to scale on \(\lambda^\perp/\lambda\), becomes
\(\beta(\alpha)\), because
\[
 (\dot y-z_0\dot v)(\dot v)=\dot y(\dot v).
\]
Let \(\Theta\subset U\oplus U^*\) be the resulting
\(r\)-dimensional image of \(T_\lambda\Lambda\).

The affine incidence over the smooth part of
\(\Lambda\cap\mathcal Q^\circ\) is an affine bundle with
fibers \(F_\lambda\), and its point image lies in \(V\) by
the preceding compactification.  At a point
\(z=z_0+D\in F_\lambda\), with
\(D\in K_v=\bigwedge^2U^*\), differentiating \(zv=y\)
gives
\[
 \dot z\,v=\dot y-z\dot v=\beta-D\alpha.
\]
The vertical tangent space maps onto \(K_v\).  Evaluation at
\(v\) identifies \(S/K_v\) with \(U^*\), so modulo these
vertical directions the differential is
\begin{equation}
 \Theta\longrightarrow U^*,
 \qquad (\alpha,\beta)\longmapsto\beta-D\alpha.
 \label{conc:deg-skew-pencil}
\end{equation}
Every parameter tangent vector lifts because the incidence is an
affine bundle.  At general points its differential has rank at
most \(\dim V=n-1\).  Hence \eqref{conc:deg-skew-pencil}
has rank less than \(r\) for general \(D\) and general
\(\lambda\).  Its determinant is polynomial in the entries of
\(D\); for each general \(\lambda\) it therefore vanishes
for every skew form \(D\) on \(U\).

\paragraph{Linear algebra of the singular skew pencil.}
We prove that this determinant condition forces \(\Theta\)
to be isotropic.  Put
\[
 B_*=\Theta\cap U^*,\qquad
 A_*=\operatorname{pr}_U(\Theta),\qquad
 H_*=B_*^\perp\subset U.
\]
If \(p=\dim A_*\), then \(\dim B_*=r-p\) and
\(\dim H_*=p\).  The plane \(\Theta\) defines a linear map
\[
 C:A_*\longrightarrow U^*/B_*\simeq H_*^*.
\]
The map in \eqref{conc:deg-skew-pencil} is the identity on
\(B_*\).  On taking the quotient by this subspace, the
determinant condition is precisely
\begin{equation}
 \det\bigl(C-D|_{A_*\times H_*}\bigr)=0
 \qquad\hbox{for every }D\in\bigwedge^2U^*.
 \label{conc:deg-quotient-pencil}
\end{equation}
Here the restriction is regarded as a map \(A_*\to H_*^*\).
If \(p=0\), the original map is the identity on \(U^*\),
a contradiction.  Thus \(p>0\).

We claim that \(A_*=H_*\).  Otherwise let
\(q=\dim(A_*\cap H_*)<p\), and choose bases
\[
 \begin{split}
 A_*&=\operatorname{span}(c_1,\ldots,c_q,a_1,\ldots,a_{p-q}),\\
 H_*&=\operatorname{span}(c_1,\ldots,c_q,h_1,\ldots,h_{p-q}).
 \end{split}
\]
All vectors displayed here are jointly linearly independent when
the common vectors are listed only once.  A restriction matrix of
an ambient skew form on \(A_*\times H_*\) may therefore be
chosen arbitrarily except that its common \(q\)-by-\(q\)
block must be skew.  Each such prescription extends to a skew
form on \(A_*+H_*\), and then to \(U\).

There is an invertible restriction matrix of this kind.  For
even \(q\), use a nondegenerate skew block on the common
indices and the identity between the remaining row and column
indices.  For odd \(q\), use a nondegenerate skew block on
\(c_1,\ldots,c_{q-1}\), the block
\[
 \begin{pmatrix}0&1\\1&0\end{pmatrix}
\]
with rows \((c_q,a_1)\) and columns \((c_q,h_1)\), and an
identity matrix on the unused extra indices.  Its two off-diagonal
entries pair different independent vectors and so can both be
chosen equal to one without violating skew symmetry.  In either
case we obtain an ambient skew form \(D_0\) with invertible
restriction.  The polynomial
\(\det(C-tD_0|_{A_*\times H_*})\) then has a nonzero leading
coefficient, contradicting \eqref{conc:deg-quotient-pencil}.
This proves \(A_*=H_*\).

If \(p\) were even, a nondegenerate skew form on \(A_*\),
extended to \(U\), would give the same contradiction.  Hence
\(p\) is odd.  We may now regard \(C\) as a bilinear form
on \(A_*\).  If it is not skew, polarization in characteristic
zero gives \(x\in A_*\) with \(C(x,x)\ne0\).
Choose a complement to \(\mathbb{C}x\), and choose a skew
form \(D_0\) which annihilates \(x\) and is nondegenerate
on this even-dimensional complement.  Extend it to \(U\).
In a basis beginning with \(x\), the coefficient of
\(t^{p-1}\) in \(\det(C-tD_0|_{A_*})\) is
\[
 (-1)^{p-1}C(x,x)
    \det\bigl(D_0|_{A_*/\mathbb{C}x}\bigr)\ne0.
\]
To obtain this power of \(t\), the determinant must use all
remaining rows and columns from \(D_0\), leaving the entry
\(C(x,x)\).  This also covers \(p=1\), with the empty
determinant equal to one.  Again we contradict
\eqref{conc:deg-quotient-pencil}; therefore \(C\) is skew.

Finally \(B_*=A_*^\perp\).  For every
\((\alpha,\beta)\in\Theta\) we consequently have
\(\beta(\alpha)=C(\alpha,\alpha)=0\).
Polarization shows that \(\Theta\) is totally isotropic.
Its dimension is \(r\) in the nondegenerate \(2r\)-dimensional
space \(U\oplus U^*\), so it is maximal isotropic.  We have
proved this tangent property on a dense open subset of the smooth
locus of \(\Lambda\).

\paragraph{Local rigidity of maximal null tangent spaces.}
We finish by showing directly that this tangent property makes
\(\Lambda\) projectively linear.  In the vector space defining
\(\mathcal Q\), choose a hyperbolic pair \(e_0,f_0\), with
orthogonal complement \(W\), and normalize the quadratic form as
\[
 q(se_0+x+tf_0)=st+q_0(x),\qquad x\in W.
\]
Choose the pair so that a general smooth point of \(\Lambda\)
lies in the chart
\begin{equation}
 W\longrightarrow\mathcal Q,
 \qquad x\longmapsto[e_0+x-q_0(x)f_0].
 \label{conc:deg-conformal-chart}
\end{equation}
The tangent quadratic form in this chart is \(q_0(dx)\):
the derivative of the displayed lift is
\(dx-(dq_0)_x(dx)f_0\), with quadratic value \(q_0(dx)\).
Shrink to a complex analytic neighborhood in the smooth locus
on which the maximal null tangent property holds everywhere.

Linear coordinates \((u,w)\in\mathbb{C}^r\oplus\mathbb{C}^r\)
on \(W\) can be chosen so that
\[
 q_0(u,w)=u\cdot u-w\cdot w,
 \qquad
 T_{x_*}\Lambda=\{(\zeta,\zeta):\zeta\in\mathbb{C}^r\}.
\]
Here the dot denotes the standard complex bilinear pairing.
For justification, choose a basis of the null tangent plane and
a dual basis in a complementary plane.  Subtracting half the
Gram matrix in the first set of directions makes the complementary
plane null as well.  Sums and differences of these bases, with
rescaling over \(\mathbb{C}\), give the displayed model.
The projection to the \(u\) coordinates is invertible on the
tangent space at \(x_*\), so the holomorphic inverse function
theorem writes the local submanifold as \(w=F(u)\).

Write \(F_i=\partial_iF\) and \(F_{ij}=\partial_i\partial_jF\).
Null tangency says
\[
 g_{ij}:=F_i\cdot F_j=\delta_{ij}.
\]
Differentiating these identities and commuting mixed partials gives
\[
 2F_{ij}\cdot F_l
   =\partial_i g_{jl}+\partial_j g_{il}-\partial_l g_{ij}=0.
\]
The vectors \(F_1,\ldots,F_r\) are a basis, since their Gram
matrix is the identity.  Thus \(F_{ij}=0\) for all \(i,j\),
and \(F\) is affine on a smaller connected neighborhood.
The local submanifold in \(W\) is an open subset of an affine
plane \(x_*+L_0\), where \(L_0\) is an \(r\)-dimensional
totally isotropic linear subspace.

Let
\(B_0(x,w)=q_0(x+w)-q_0(x)-q_0(w)\) be the polarization.
For \(w\in L_0\),
\(q_0(x_*+w)=q_0(x_*)+B_0(x_*,w)\), so the lift in
\eqref{conc:deg-conformal-chart} is affine-linear on this plane.
Its projective closure is the projectivization of
\[
 \mathbb{C}\bigl(e_0+x_*-q_0(x_*)f_0\bigr)
 \; +\;
 \{w-B_0(x_*,w)f_0:w\in L_0\}.
\]
This vector space has dimension \(r+1\): the first summand
has nonzero \(e_0\) coordinate, and the second maps injectively
to \(L_0\).  Its quadratic form vanishes identically.  Indeed,
it vanishes on the displayed affine lift, and homogeneity extends
this identity to its linear span.  We have obtained a projective
linear \(\mathbf{P}^r\subset\mathcal Q\) containing a
nonempty analytic open subset of the smooth irreducible variety
\(\Lambda\).

Such an open subset is Zariski dense, since a proper algebraic
subset has smaller dimension and cannot contain a smooth
neighborhood.  Hence \(\Lambda\) is contained in this
\(\mathbf{P}^r\), and equality follows from
\(\dim\Lambda=r\).  This proves the exceptional assertion
and completes the proof.
\end{proof}


\subsection{Closing the estimates with uniform constants}
\label{conc:uniform-completion}

We now prove Theorem~\ref{conc:theorem}.  The preceding results control
flats away from a low-degree exceptional section when $V$ has a
low-degree approximate description.  A description containing a large
degree instead produces a larger carrier whose degree is sufficiently
smaller.  A finite system of inequalities combines the two alternatives.

\begin{proof}[Proof of Theorem~\ref{conc:theorem}]
Fix the sample, $A$, and one of Specifications (a) and (b).  For
$0\leq h<r$, let $R_h$ be the supremum of
\[
 \frac{\#\{\lambda\in\mathcal P:F_\lambda\subset V\}}
      {\deg(V)A^h}
\]
over the allowed irreducible varieties of dimension $s_0+h$.
An empty supremum is zero.  These suprema are finite for the fixed
sample: each numerator is at most $|\mathcal P|$, and each denominator
is at least one.  As noted after the theorem, $R_0\leq1$.  In (a) also
put
\[
 R_r=|\mathcal P|/A^r\leq C_0.
\]
In (b) the index $r$ is unavailable and will never be used.

We first refine the choice of approximate equations.  Take the
$c=n-\dim V$ equations from Theorem~\ref{conc:walsh}, ordering their
degrees as $d_1\leq\cdots\leq d_c$.  They form a parameter sequence
in the regular local ring $\mathbb C[S]_{I(V)}$, so every initial
segment of length $i$ has height $i$
\mbox{\cite[Tags~00NQ, 02JN, and 00NA]{stacksproject}}.  Choose a polynomial $q$ of
smallest positive degree in $I(V)$.  We may use $q$ as the first
equation and preserve the upper bounds $d_i$ on the remaining ones.
Indeed, suppose a replacement prefix of length $i-1$ has height
$i-1$.  Its minimal primes have that height, because a partial
parameter sequence in this regular local ring is a regular sequence.
None can contain the ideal generated by the first $i$ original
equations, whose height is $i$.  A generic constant linear combination
of those $i$ equations avoids all these minimal primes and extends
the replacement prefix to height $i$.  Its degree is at most $d_i$.
Starting with $q$ and continuing gives a parameter system again.
Consequently $V$ remains a component of its global common zero set.
Replace $d_1$ by $\deg q$.  The nondecreasing upper bounds still satisfy
\begin{equation}
 \prod_{i=1}^c d_i\leq C_n\deg V.
 \label{conc:description-product}
\end{equation}
The later actual degrees may be smaller than their upper bounds;
this only strengthens the relevant B\'ezout estimates.

Choose constants
\[
 K_1>K_2>\cdots>K_{r-1}>C_0,
\]
to be fixed below in terms of $e,C_0$.  Consider an allowed carrier
$V$ of dimension $s_0+h$, with $1\leq h<r$.

Suppose first that every description degree $d_i$ is at most $K_hA$.
Lemma~\ref{conc:rulings} supplies a proper exceptional section
$V\cap Z(H)$ with
\[
 \deg H\leq C_eK_hA.
\]
Outside it, there are boundedly many ruling incidences.  For one with
$j=0$, Proposition~\ref{conc:degree} gives a parameter variety of
dimension $h$ and degree at most $C_e\deg V$; the separator bound
then counts at most $C_{e,C_0}\deg(V)A^h$ parameters on it.  For one
with $j=1$, that proposition says $h=r-1$ and the parameter variety is
a maximal isotropic $\mathbf P^r$.  Assumption (ii) counts at most
$C_0A^h$ parameters on it.  Thus all represented flats together cost
at most $C_{e,C_0}\deg(V)A^h$.

The remaining flats are in $V\cap Z(H)$.  Its components, if any,
have dimension $s_0+h-1$ and total degree at most
$C_eK_hA\deg V$.  They remain allowed in (b), since they are subsets
of the original allowed carrier.  The definition of $R_{h-1}$ gives
\begin{equation}
 \frac{\#\{\lambda:F_\lambda\subset V\}}{\deg(V)A^h}
 \leq C(1+K_hR_{h-1}).
 \label{conc:low-recurrence}
\end{equation}
Here and below $C$ depends only on $e,C_0$, not on the $K_i$.

Suppose instead that some $d_i>K_hA$.  Let $p<c$ be the length of
the prefix with $d_i\leq K_hA$.  The local prefix is a regular
sequence of height $p$.  A minimal prime of it in
$\mathbb C[S]_{I(V)}$ contracts to a global component $W$ of that
prefix, containing $V$ and of codimension $p$.  Isolated-component
B\'ezout gives
\[
 \deg W\leq\prod_{i\leq p}d_i.
\]
For $p=0$, take $W=S$ and the empty product to be one.  Put
$i=r-p>h$, so $\dim W=s_0+i$.  From
\eqref{conc:description-product},
\begin{equation}
 \deg V\geq C_n^{-1}\deg W(K_hA)^{i-h}.
 \label{conc:degree-comparison}
\end{equation}
In Case (b), the minimal first equation has degree at most $C_0A$.
Since $K_h>C_0$, it belongs to the prefix.  Therefore $p\geq1$,
$i<r$, and $W$ lies in that same hypersurface of degree at most
$C_0A$.  Thus $W$ is allowed in (b).  In (a) the ambient possibility
$i=r$ is already controlled.  In either case, counting first on $W$
yields
\begin{equation}
 \frac{\#\{\lambda:F_\lambda\subset V\}}{\deg(V)A^h}
 \leq \frac{C R_i}{K_h^{i-h}}.
 \label{conc:high-recurrence}
\end{equation}

Taking suprema in the two alternatives gives the finite system
\begin{equation}
 R_h\leq\max\left\{
 C(1+K_hR_{h-1}),\quad
 \max_{\substack{i>h\\i\text{ available}}}
       \frac{CR_i}{K_h^{i-h}}\right\},
 \qquad1\leq h<r.
 \label{conc:recurrence-system}
\end{equation}
An empty inner maximum is omitted.  We close this system explicitly,
so that no constant depends on the description degrees or on the
size of the sample.

Enlarge $C$ to at least one.  Choose
\[
 B_0\geq\max(1,C_0),\qquad c_1>4C+4,\qquad \kappa>4Cc_1.
\]
Take $K_{r-1}\geq\max(2,C_0+1)$, and successively set
$K_h=\kappa K_{h+1}$ for $h=r-2,\ldots,1$.  Define barriers
\[
 B_h=c_1K_hB_{h-1}\quad(1\leq h<r),
 \qquad B_r=B_0\quad\hbox{in Case (a)}.
\]
For interior indices $h<i<r$,
\[
 \frac{CB_i}{B_hK_h^{i-h}}
 =C c_1^{i-h}\prod_{j=h+1}^{i}\frac{K_j}{K_h}
 \leq C(c_1/\kappa)^{i-h}<\frac14.
\]
If the ambient index is available, $B_h\geq c_1K_hB_0$ gives
\[
 \frac{CB_r}{B_hK_h^{r-h}}
 \leq\frac{C}{c_1K_h^{r-h+1}}<\frac12.
\]
The other transition satisfies
\[
 \frac{C(1+K_hB_{h-1})}{B_h}\leq\frac{2C}{c_1}<\frac12.
\]
All choices so far depend only on $e,C_0$.

If the maximum $T=\max_h R_h/B_h$ were greater than one, a maximizing
index would be interior: $R_0\leq B_0$, and the ambient endpoint,
when present, obeys the same bound.  Substitute $R_j\leq TB_j$ into
\eqref{conc:recurrence-system}.  Since $T>1$, each term in its right
side is strictly less than $TB_h$, contradicting maximality.  Hence
$R_h\leq B_h$ for every available index.  This is
\eqref{conc:bound}, with a constant depending only on $e,C_0$.
\end{proof}
